\chapter{Methodology}

This chapter presents the proposed Evidential Fourier Neural Operator (Evidential-FNO), a novel approach that integrates Evidential Deep Learning with Fourier Neural Operators for uncertainty-aware PDE solution prediction. We first establish the theoretical connection between FNO outputs and regression problems, then describe the architectural modifications required to enable evidential learning.

\section{Evidential Fourier Neural Operator}

\subsection{FNO Output as Spatially-Extended Regression}

The key insight enabling the integration of Evidential Deep Learning with Fourier Neural Operators lies in recognizing the nature of FNO's output. Consider the standard FNO that learns a solution operator $\mathcal{G}_\theta: \mathcal{A} \rightarrow \mathcal{U}$ mapping input functions to solution functions. For a discretized input $a \in \mathbb{R}^{H \times W \times d_{in}}$ on a spatial grid, the vanilla FNO produces:
\begin{equation}
    u = \mathcal{G}_\theta(a) \in \mathbb{R}^{H \times W \times 1}
\end{equation}
where each spatial location $(i, j)$ yields a continuous scalar prediction $u_{i,j} \in \mathbb{R}$.

This output structure is fundamentally equivalent to performing regression at each spatial location simultaneously. At any grid point $(i, j)$, the FNO effectively solves a regression problem:
\begin{equation}
    u_{i,j} = f_\theta(a, i, j) + \epsilon_{i,j}, \quad \epsilon_{i,j} \sim \mathcal{N}(0, \sigma^2_{i,j})
\end{equation}
where $f_\theta$ represents the learned mapping and $\epsilon_{i,j}$ captures prediction noise that may vary spatially.

The crucial observation is that while the FNO operates on 2D spatial grids, each output value remains a continuous scalar---precisely the setting where Evidential Deep Learning excels. Standard regression neural networks predict point estimates; Evidential Deep Learning extends this by predicting distributions over possible outputs. This same extension applies naturally to the spatially-extended regression performed by FNOs.

\subsection{From Point Estimates to Evidential Outputs}

In standard regression, a neural network predicts a single value $\hat{y}$ for each input. Evidential regression instead predicts parameters of a higher-order distribution---specifically, a Normal-Inverse-Gamma (NIG) distribution that represents uncertainty over both the prediction mean and variance.

For the FNO setting, we extend this from scalar regression to spatially-extended regression. Rather than predicting a single channel output $u \in \mathbb{R}^{H \times W \times 1}$, the Evidential-FNO predicts four channels representing the NIG parameters at each spatial location:
\begin{equation}
    \mathcal{G}^{\text{evid}}_\theta(a) = \begin{pmatrix} \gamma \\ \nu \\ \alpha \\ \beta \end{pmatrix} \in \mathbb{R}^{H \times W \times 4}
\end{equation}

At each spatial location $(i,j)$, the four parameters define a NIG distribution:
\begin{equation}
    (\mu_{i,j}, \sigma^2_{i,j}) \sim \text{NIG}(\gamma_{i,j}, \nu_{i,j}, \alpha_{i,j}, \beta_{i,j})
\end{equation}
which serves as a prior over the Gaussian likelihood parameters $(\mu, \sigma^2)$ for that location's prediction.

The NIG distribution factorizes as:
\begin{align}
    \sigma^2_{i,j} &\sim \text{Inv-Gamma}(\alpha_{i,j}, \beta_{i,j}) \\
    \mu_{i,j} | \sigma^2_{i,j} &\sim \mathcal{N}\left(\gamma_{i,j}, \frac{\sigma^2_{i,j}}{\nu_{i,j}}\right)
\end{align}

This formulation provides closed-form expressions for predictive statistics at each grid point:
\begin{align}
    \mathbb{E}[u_{i,j}] &= \gamma_{i,j} \label{eq:pred_mean} \\
    \text{Var}[u_{i,j}] &= \frac{\beta_{i,j}}{\alpha_{i,j} - 1} \cdot \frac{\nu_{i,j} + 1}{\nu_{i,j}} \label{eq:pred_var}
\end{align}

\subsection{Uncertainty Decomposition}

A key advantage of the evidential formulation is the natural decomposition of predictive uncertainty into epistemic and aleatoric components. For each spatial location $(i,j)$:

\textbf{Aleatoric uncertainty} (irreducible data noise):
\begin{equation}
    \sigma^2_{\text{ale}, i,j} = \mathbb{E}[\sigma^2_{i,j}] = \frac{\beta_{i,j}}{\alpha_{i,j} - 1}
\end{equation}

\textbf{Epistemic uncertainty} (model uncertainty):
\begin{equation}
    \sigma^2_{\text{epi}, i,j} = \text{Var}[\mu_{i,j}] = \frac{\beta_{i,j}}{\nu_{i,j}(\alpha_{i,j} - 1)}
\end{equation}

The epistemic uncertainty is inversely proportional to $\nu$, which represents the ``evidence'' or effective number of virtual observations supporting the prediction. Regions where the model has seen similar training data will have higher $\nu$ and thus lower epistemic uncertainty.

\subsection{Architectural Modification}

The architectural change from vanilla FNO to Evidential-FNO is minimal, affecting only the final projection layer. The complete architecture is:

\textbf{Vanilla FNO:}
\begin{equation}
    a \xrightarrow{P} v_0 \xrightarrow{\mathcal{F}_1} v_1 \xrightarrow{\mathcal{F}_2} v_2 \xrightarrow{\mathcal{F}_3} v_3 \xrightarrow{\mathcal{F}_4} v_4 \xrightarrow{Q} u \in \mathbb{R}^{H \times W \times 1}
\end{equation}

\textbf{Evidential-FNO:}
\begin{equation}
    a \xrightarrow{P} v_0 \xrightarrow{\mathcal{F}_1} v_1 \xrightarrow{\mathcal{F}_2} v_2 \xrightarrow{\mathcal{F}_3} v_3 \xrightarrow{\mathcal{F}_4} v_4 \xrightarrow{Q^{\text{evid}}} (\gamma, \nu, \alpha, \beta) \in \mathbb{R}^{H \times W \times 4}
\end{equation}

where:
\begin{itemize}
    \item $P: \mathbb{R}^{d_{in}} \rightarrow \mathbb{R}^{w}$ is the lifting layer (unchanged)
    \item $\mathcal{F}_l$ are spectral convolution layers (unchanged)
    \item $Q: \mathbb{R}^{w} \rightarrow \mathbb{R}^{1}$ is the vanilla projection
    \item $Q^{\text{evid}}: \mathbb{R}^{w} \rightarrow \mathbb{R}^{4}$ is the evidential projection
\end{itemize}

The evidential projection layer applies appropriate activation functions to ensure valid NIG parameters:
\begin{align}
    \gamma &= \text{linear}(v_4) && \text{(unbounded, prediction mean)} \\
    \nu &= \text{softplus}(v_4) + \epsilon && \text{(positive, evidence)} \\
    \alpha &= \text{softplus}(v_4) + 1 + \epsilon && \text{(> 1, shape parameter)} \\
    \beta &= \text{softplus}(v_4) + \epsilon && \text{(positive, scale parameter)}
\end{align}
where $\epsilon = 10^{-6}$ ensures numerical stability.

\subsection{Training Objective}

The Evidential-FNO is trained by minimizing the negative log-likelihood of the NIG distribution. For a training pair $(a, u^*)$ where $u^* \in \mathbb{R}^{H \times W}$ is the ground truth solution:
\begin{equation}
    \mathcal{L}_{\text{NIG}} = \frac{1}{HW} \sum_{i,j} \mathcal{L}^{(i,j)}_{\text{NIG}}
\end{equation}
where the per-location loss is:
\begin{equation}
    \mathcal{L}^{(i,j)}_{\text{NIG}} = \frac{1}{2}\log\left(\frac{\pi}{\nu_{i,j}}\right) - \alpha_{i,j}\log(\Omega_{i,j}) + \left(\alpha_{i,j} + \frac{1}{2}\right)\log\left((u^*_{i,j} - \gamma_{i,j})^2 \nu_{i,j} + \Omega_{i,j}\right) + \log\frac{\Gamma(\alpha_{i,j})}{\Gamma(\alpha_{i,j} + \frac{1}{2})}
\end{equation}
with $\Omega_{i,j} = 2\beta_{i,j}(1 + \nu_{i,j})$.

\subsection{Regularization Schemes}

A known challenge in evidential deep learning is preventing the model from ``explaining away'' prediction errors by inflating uncertainty estimates rather than improving predictions. This occurs when the model increases $\beta$ (aleatoric uncertainty) to reduce the NIG loss without actually learning better predictions. Regularization is essential to encourage meaningful uncertainty estimates.

We investigate six regularization schemes, each with different inductive biases for controlling the evidence parameters:

\subsubsection{Standard Regularization}

The original evidential regularization from \cite{amini2020deep} penalizes evidence on incorrect predictions:
\begin{equation}
    \mathcal{L}_{\text{standard}} = \mathcal{L}_{\text{NIG}} + \lambda \cdot |u^* - \gamma| \cdot (2\nu + \alpha)
\end{equation}

\textbf{Intuition}: When prediction error $|u^* - \gamma|$ is high, the model is penalized for having high evidence ($\nu$) and concentration ($\alpha$). This encourages the model to express high confidence only when predictions are accurate.

\subsubsection{Improved Regularization}

We propose an improved scheme that directly regularizes the evidence parameter toward a baseline value:
\begin{equation}
    \mathcal{L}_{\text{improved}} = \mathcal{L}_{\text{NIG}} + \lambda \cdot |u^* - \gamma| \cdot (2\nu + \alpha) + \lambda_{\nu} \cdot |\nu - 1|
\end{equation}

\textbf{Intuition}: The additional term $|\nu - 1|$ prevents $\nu$ from growing unboundedly large or collapsing to near-zero. By anchoring $\nu$ toward 1, we maintain a balance between over-confident and under-confident predictions. The evidence can still deviate from 1 when the data supports it, but extreme values are discouraged.

\subsubsection{Uncertainty-Aware Regularization}

This scheme adaptively scales regularization based on prediction error magnitude:
\begin{equation}
    \mathcal{L}_{\text{UA}} = \mathcal{L}_{\text{NIG}} + \lambda \cdot (u^* - \gamma)^2 \cdot \nu
\end{equation}

\textbf{Intuition}: The quadratic error term $(u^* - \gamma)^2$ creates stronger penalties for large errors. Multiplying by $\nu$ means that high-evidence predictions are especially penalized when wrong. This creates a self-correcting mechanism: the model learns to reduce evidence in regions where it makes large errors.

\subsubsection{Adaptive Regularization}

We introduce epoch-dependent regularization strength that evolves during training:
\begin{equation}
    \mathcal{L}_{\text{adaptive}} = \mathcal{L}_{\text{NIG}} + \lambda(t) \cdot |u^* - \gamma| \cdot \phi(\nu, \alpha)
\end{equation}
where:
\begin{equation}
    \lambda(t) = \lambda_0 \cdot \min\left(1, \frac{t}{t_{\text{warmup}}}\right)
\end{equation}
and $\phi(\nu, \alpha) = 2\nu + \alpha$ is the evidence function.

\textbf{Intuition}: Early in training, predictions are poor, so heavy regularization would prevent the model from learning. The warmup schedule starts with weak regularization ($\lambda \approx 0$), allowing the model to first learn accurate predictions. As training progresses, regularization strengthens to refine uncertainty estimates. This two-phase approach separates ``learning to predict'' from ``learning to quantify uncertainty.''

\subsubsection{L2-Evidence Regularization}

This scheme applies L2 (squared) penalty directly on the evidence parameter:
\begin{equation}
    \mathcal{L}_{\text{L2}} = \mathcal{L}_{\text{NIG}} + \lambda \cdot |u^* - \gamma| \cdot (2\nu + \alpha) + \lambda_{\nu} \cdot \nu^2
\end{equation}

\textbf{Intuition}: Unlike L1 regularization which encourages sparsity, L2 regularization on $\nu^2$ creates smooth, distributed evidence values. Large $\nu$ values are penalized quadratically, strongly discouraging overconfident predictions. This is analogous to weight decay in standard neural network training---it prevents any single prediction from having disproportionately high evidence.

\subsubsection{KL-Divergence Regularization}

This scheme penalizes deviation from a prior NIG distribution:
\begin{equation}
    \mathcal{L}_{\text{KL}} = \mathcal{L}_{\text{NIG}} + \lambda \cdot |u^* - \gamma| \cdot (2\nu + \alpha) + \lambda_{\text{KL}} \cdot D_{\text{KL}}\left(\text{NIG}(\gamma, \nu, \alpha, \beta) \| \text{NIG}(\gamma_0, \nu_0, \alpha_0, \beta_0)\right)
\end{equation}

The KL divergence between NIG distributions has closed form:
\begin{equation}
    D_{\text{KL}} = \frac{1}{2}\left[\frac{\nu_0}{\nu} - \log\frac{\nu_0}{\nu} - 1\right] + \alpha_0\log\frac{\beta}{\beta_0} - \log\frac{\Gamma(\alpha)}{\Gamma(\alpha_0)} + (\alpha - \alpha_0)\psi(\alpha) + \frac{\beta_0 - \beta}{\beta}\alpha
\end{equation}
where $\psi(\cdot)$ is the digamma function and $(\gamma_0, \nu_0, \alpha_0, \beta_0)$ are prior hyperparameters.

\textbf{Intuition}: This is the most principled Bayesian approach. The KL term acts as a prior, pulling the predicted distribution toward a ``default'' uncertainty level. Predictions can deviate from the prior when evidence supports it, but the model must ``pay'' an information-theoretic cost. This naturally implements Occam's razor for uncertainty estimation.

\subsubsection{Summary of Regularization Schemes}

Table~\ref{tab:reg_summary} summarizes the key characteristics of each regularization scheme:

\begin{table}[h]
\centering
\caption{Comparison of regularization schemes for Evidential-FNO.}
\label{tab:reg_summary}
\begin{tabular}{lll}
\toprule
\textbf{Scheme} & \textbf{Key Mechanism} & \textbf{Target Behavior} \\
\midrule
Standard & Error-weighted evidence penalty & Reduce confidence on errors \\
Improved & Evidence anchoring to baseline & Prevent extreme $\nu$ values \\
Uncertainty-Aware & Quadratic error scaling & Strong penalty for large errors \\
Adaptive & Epoch-dependent strength & Learn predictions before uncertainty \\
L2-Evidence & Squared evidence penalty & Smooth, distributed confidence \\
KL-Divergence & Prior distribution matching & Bayesian uncertainty regularization \\
\bottomrule
\end{tabular}
\end{table}

The total loss for all schemes follows the general form:
\begin{equation}
    \mathcal{L} = \mathcal{L}_{\text{NIG}} + \mathcal{L}_{\text{reg}}
\end{equation}
where $\mathcal{L}_{\text{reg}}$ varies according to the chosen scheme. All schemes share the same network architecture; only the training objective differs.

\subsection{Computational Efficiency}

A significant advantage of Evidential-FNO over ensemble methods is computational efficiency. Table~\ref{tab:computational} summarizes the comparison:

\begin{table}[h]
\centering
\caption{Computational comparison of uncertainty quantification methods for FNO.}
\label{tab:computational}
\begin{tabular}{lccc}
\toprule
\textbf{Method} & \textbf{Training Cost} & \textbf{Inference Cost} & \textbf{Parameters} \\
\midrule
Vanilla FNO & $1\times$ & $1\times$ & $1\times$ \\
MC Dropout & $1\times$ & $T\times$ & $1\times$ \\
Deep Ensemble ($M$ models) & $M\times$ & $M\times$ & $M\times$ \\
Evidential-FNO & $1\times$ & $1\times$ & $\approx 1\times$ \\
\bottomrule
\end{tabular}
\end{table}

The Evidential-FNO requires only a single forward pass during both training and inference, with negligible parameter overhead (only the final projection layer is modified). This makes it particularly attractive for real-time applications where computational budget is limited.

\section{Baseline Methods}

For comprehensive evaluation, we compare Evidential-FNO against established uncertainty quantification methods adapted for neural operators.

\subsection{MC Dropout}

Monte Carlo Dropout approximates Bayesian inference by using dropout at test time. We apply Dropout2d (spatial dropout) after each Fourier layer with dropout rate $p$. At inference, $T$ stochastic forward passes yield:
\begin{align}
    \hat{u}_{i,j} &= \frac{1}{T}\sum_{t=1}^{T} u^{(t)}_{i,j} \\
    \hat{\sigma}^2_{i,j} &= \frac{1}{T}\sum_{t=1}^{T} (u^{(t)}_{i,j} - \hat{u}_{i,j})^2
\end{align}

\subsection{Deep Ensemble}

Deep Ensemble trains $M$ independent FNO models with different random initializations. Predictions are aggregated as:
\begin{align}
    \hat{u}_{i,j} &= \frac{1}{M}\sum_{m=1}^{M} u^{(m)}_{i,j} \\
    \hat{\sigma}^2_{i,j} &= \frac{1}{M}\sum_{m=1}^{M} (u^{(m)}_{i,j} - \hat{u}_{i,j})^2
\end{align}

\section{Summary}

This chapter presented Evidential-FNO, which extends Fourier Neural Operators with uncertainty quantification through Evidential Deep Learning. The key contributions are:

\begin{enumerate}
    \item \textbf{Theoretical connection}: We established that FNO's continuous 2D output is equivalent to spatially-extended regression, enabling direct application of evidential learning.

    \item \textbf{Architectural modification}: The modification is minimal---only the output projection changes from 1 to 4 channels, with appropriate activation functions ensuring valid NIG parameters.

    \item \textbf{Uncertainty decomposition}: The evidential formulation provides principled decomposition into epistemic and aleatoric uncertainty at each spatial location.

    \item \textbf{Regularization schemes}: We investigate six regularization strategies---Standard, Improved, Uncertainty-Aware, Adaptive, L2-Evidence, and KL-Divergence---each with different inductive biases for controlling evidence estimation.

    \item \textbf{Computational efficiency}: Unlike ensemble or dropout methods, Evidential-FNO requires only a single forward pass for uncertainty estimation.
\end{enumerate}